% CV classique une colonne. Option fond: \usepackage{pagecolor} \pagecolor{gray!8}
\documentclass[11pt,a4paper]{article}
\usepackage{fontspec}
\usepackage{graphicx}
\usepackage{hyperref}
\usepackage{geometry}
\geometry{left=1.5cm,right=1.5cm,top=1.5cm,bottom=1.5cm}
\IfFontExistsTF{TeX Gyre Termes}{\setmainfont{TeX Gyre Termes}}{\IfFontExistsTF{Liberation Serif}{\setmainfont{Liberation Serif}}{\setmainfont{Times New Roman}}}

\title{\vspace{-1.2cm}\textbf{Votre nom}}
\author{}
\date{}

\begin{document}
\maketitle
\vspace{-0.8cm}
% Photo : remplacer \rule par \includegraphics[width=2.2cm,clip]{photo}
\begin{flushright}\rule{2.2cm}{2.8cm}\end{flushright}

\textbf{Contact :} email@example.com \quad | \quad Téléphone \quad | \quad Ville, Pays \\
\textbf{LinkedIn :} linkedin.com/in/yourprofile

\vspace{0.4cm}
\noindent\rule{\textwidth}{0.4pt}

\section*{Formation}
\textbf{Diplôme} \hfill Année -- Année \\
Établissement, Ville \\
Brève mention (honneurs, notes).

\section*{Expérience}
\textbf{Poste} \hfill Mois Année -- Présent \\
Entreprise, Ville \\
\begin{itemize}
  \item Réalisation ou responsabilité.
  \item Autre point.
\end{itemize}

\textbf{Poste précédent} \hfill Mois Année -- Mois Année \\
Entreprise \\
\begin{itemize}
  \item Responsabilité ou résultat clé.
\end{itemize}

\section*{Compétences}
Techniques : Compétence 1, 2, 3. \\
Langues : Langue 1 (niveau), Langue 2 (niveau).

\end{document}
